\section{Benchmark Design and Experimental Setup}
\label{sec:methods}

\subsection{Template Grammar and Complexity Tiers}
\label{sec:methods-templates}
DyMuT is built from sixty hand-authored templates, each specifying a body skeleton with parameter placeholders, a parameter-domain specification, a reference oracle implementation, and a docstring template with synonym variation. A \emph{benchmark instantiation} is one full generator run: fifteen templates per tier, four tiers, five instances per template, giving 300 problems per instantiation. Templates are stratified across four cyclomatic-complexity (CC) tiers~\cite{mccabe1976complexity} -- L1 (Trivial, CC 1, single-path pure transformations), L2 (Branching, CC 2--3, disjoint conditional logic), L3 (Iterative, CC 4--7, one or two loops with boundary conditions), and L4 (Compositional, CC ${\geq}8$, nested branches and loops with multiple failure modes), with target mutant-count bands of 5-10, 10-15, 15-25, and 25-40 respectively. Cyclomatic complexity is computed statically at template-validation time and recomputed on every generated instance to confirm tier membership, since parameter sampling can shift a generated instance's control-flow shape. Prior to use, the full sixty-template corpus passed a within-tier diversity check and a tier-separation sanity check, both required by the study's corpus validation gate.

\subsection{Seeded Instantiation and Oracle-by-Construction Scoring}
\label{sec:methods-instantiation}
Each template yields fresh problem instances by sampling its parameter domain: variable names, literal values, constants, and docstring phrasing all vary across instances, under a deterministic seed, so a given (template, seed) pair is reproduced identically across runs while distinct seeds produce different problem content. This is what separates DyMuT from a static release. Its problem space is defined by the template grammar and the parameter domains, not by a fixed set of instances. No single generated instantiation can be memorised by a future model. Correctness comes from construction. Every generated instance carries a reference implementation, produced alongside it from the same template, so scoring never depends on a separately curated set of expected outputs.

\subsection{Mutation Operators and Metrics}
\label{sec:methods-metrics}
Test quality is scored by mutation testing~\cite{demillo1978hints, jia2011mutation}, whose mutation score is an established proxy for a suite's fault-detection ability~\cite{andrews2005mutation}, using seven semantic-altering-change (SAC) operators: Arithmetic Operator Replacement, Relational Operator Replacement, Logical Operator Replacement, Constant Replacement, Boundary Mutation, Statement Deletion, and Return/Statement Mutation, applied as single-edit mutants. Semantic-preserving-change operators are out of scope for this study. Four metrics are recorded per problem: \emph{Exec@K}, whether the generated suite is imported and collected by the test runner without error (a gate, not a quality measure); \emph{baseline\_pass\_rate}, whether the suite passes against the original, correct implementation (the precondition for mutation scoring); \emph{Mut@K}, the fraction of mutants killed; and \emph{coverage\_adjusted}, Mut@K multiplied by baseline\_pass\_rate. Naming follows Pass@K convention, where $K$ denotes attempts per problem; every result in this study uses $K=1$. The number of sampled instances per template is denoted by $m$. $N$ denotes the number of independent benchmark instantiations (full regenerations under a fresh seed).

Mut@K is \emph{conditional}: its mean is taken only over problems whose generated suite already passes the baseline-correctness gate, excluding non-passing problems from the average entirely rather than scoring them zero. This is the pre-registered primary metric throughout this study. \texttt{coverage\_adjusted} is an unconditional reading of the same data, identical to averaging Mut@K over every attempted problem with non-passing suites (scored zero) and is reported as a labelled co-primary, not a replacement. The two readings measure different quantities and can disagree widely for a model with a low baseline-pass rate.

\subsection{Models}
\label{sec:methods-models}
Nine models spanning four providers were evaluated: three closed-source frontier/efficient pairs from OpenAI (\texttt{gpt-5}/\texttt{gpt-5-mini}), from Anthropic (Claude Sonnet 4.6/Claude Haiku 4.5), and from Google (Gemini 3.1 Pro Preview/Flash-Lite). A three-point open-source scaling sweep (Qwen2.5-Coder~\cite{hui2024qwen25coder}, 32B/7B/1.5B parameters). The Google pair is a forced substitution for the originally specified Gemini 2.5 Pro/Flash, which returned model-unavailable errors ahead of their provider's own stated retirement schedule. Every model call is logged with model identifier, snapshot or revision, decoding parameters, seed, and timestamp; fine-tuned checkpoints are treated as a distinct model identity from their base weights, since adapter identity changes what is being measured even when the base model field is unchanged. % cite:external -- "1.5"/"3.1"/"2.5" are model name/version/parameter-count identifiers (Qwen2.5-Coder, Gemini 3.1, Gemini 2.5), not result values

\subsection{Fine-Tuning on ULT to Create Contamination}
\label{sec:methods-finetune}
To identify contamination-driven score inflation causally, not just observationally, the smallest open-source model (Qwen2.5-Coder-1.5B-Instruct) is fine-tuned via low-rank adaptation (LoRA)~\cite{hu2022lora} directly on ULT's ground-truth test suites, restricted to the subset of ULT problems for which a repair-verified training pair could be constructed (112 of the 400-problem evaluation subsample; the remaining 288 problems, structurally unrepairable by the same verification pipeline, form a held-out unseen-ULT population). The fine-tuned checkpoint is then re-evaluated on ULT, on PLT, and on a fresh DyMuT instantiation, using a $2{\times}3$ difference-in-differences design~\cite{angrist2009mostly} (pre/post fine-tuning $\times$ \{ULT, PLT, dynamic\}) with bootstrap confidence intervals~\cite{efron1993bootstrap}, following this study's preference for effect sizes with intervals over point-estimate hypothesis tests. A DyMuT control population restricted to the two lower complexity tiers (L1--L3, outside ULT/PLT's own complexity floor) separates any generic fine-tuning improvement from contamination-specific inflation. % cite:external -- "1.5" is the model's own parameter-count identifier (Qwen2.5-Coder-1.5B-Instruct), not a result value

\subsection{Equivalent-Mutant Validation}
\label{sec:methods-equivmut}
The mutation operators applied here can generate mutants that are semantically equivalent to the original program and therefore unkillable by any correct test suite, which would depress the raw mutation score without reflecting a genuine test-quality gap. The reported score throughout this study is the raw, unadjusted mutation score, treated as a conservative and comparable headline figure. A validated companion score is separately obtained via manual inspection of a stratified sample of 120 mutants (thirty per complexity tier), from which an estimated equivalent-mutant rate with confidence intervals is derived and applied as a correction to the sampled population. \textbf{This inspection is manual, hand-inspected work not yet complete at the time of writing}: the raw, unadjusted mutation score is reported as the default headline figure throughout this paper, and the validated companion score will be reported alongside it once the inspection sample is complete.
